\documentclass[11pt,a4paper]{article}
\usepackage[margin=2.2cm]{geometry}
\usepackage{booktabs}
\usepackage{parskip}
\usepackage[hidelinks]{hyperref}
\pagestyle{plain}
\begin{document}
\begin{center}
{\LARGE\bfseries The Odyssey}\\[4pt]
{\large A running log of the Safety-Dial project, one entry per task.}\\[2pt]
Sachin Mohan, supervised by Geraud Nangue Tasse.
\end{center}

\section*{Task 1: Replicate the ManiGuard Jar benchmark}

\textbf{Why.} Before trying to improve the Jar policy we needed to know its real base rate, and whether our setup even matches the paper. The old runs never could: they used a different render profile (640-pixel high fidelity instead of the paper's 256-pixel legacy) on two tasks.

\textbf{What we did.}
\begin{itemize}\setlength\itemsep{1pt}
  \item Pinned ManiGuard to the last revision that still renders the paper's profile, and ran the paper's protocol exactly: 26 Jar tasks, 2{,}000 steps, 8 of 16 actions, official evaluator.
  \item Rented one two-GPU machine on Vast and ran seeds 1 and 2, then seed 0 twice (one per GPU) to see how much a repeated seed moves.
  \item Dropped an earlier seed 0 run from a different machine so that everything of record comes from one host. With it the test still passed; without it the safe-success rate goes from 15\% to 23\%.
  \item Saved every rollout video to a private Hugging Face dataset so anyone can watch what the robot did. Total cost about \$6.
\end{itemize}

\textbf{Result.} Counts over 26 tasks per run. The paper's row is for comparison.

\begin{center}
\begin{tabular}{lrrrrr}
\toprule
 & Safe success & Unsafe success & Safe fail & Unsafe fail & Violations \\
\midrule
Seed 0, repeat 1 & 7 & 0 & 9 & 10 & 10 \\
Seed 0, repeat 2 & 7 & 5 & 7 & 7 & 12 \\
Seed 1 & 6 & 0 & 10 & 10 & 10 \\
Seed 2 & 4 & 2 & 11 & 9 & 11 \\
\midrule
Ours, 104 episodes & 24 (23\%) & 7 & 37 & 36 & 43 (41\%) \\
Paper, 78 episodes & 21 (27\%) & 4 & 27 & 26 & 30 (38\%) \\
\bottomrule
\end{tabular}
\end{center}

Fisher tests against the paper: $p = 0.60$ for safe success, $p = 0.76$ for violations. Our setup matches the paper.

\textbf{What we learned.}
\begin{itemize}\setlength\itemsep{1pt}
  \item Running the same seed twice gives the same outcome on only 13 of 26 tasks, about the same as two different seeds. The robot reaches the jar the same way every time and then the physics at contact decides. One episode is a coin flip with a per-task bias.
  \item 13 tasks never succeed safely, 12 sit in the middle, one always succeeds. The middle 12 are where a dial could show anything.
  \item Most violations are the jar tipping over (33 of 43); the gripper only closed on the jar in 32 of 104 episodes.
  \item To tell a ten-point improvement from noise we will need about 200 episodes per condition. Everything after this is sized on that.
\end{itemize}

\textbf{Where it lives.} Videos and rows: \texttt{Sachioster/maniguard-jar-replication-20261022} (private). Receipts: \texttt{progress/receipts/}. Full write-up: \emph{Jar replication and noise floor} under \texttt{docs/research/research\_notes/}.
\newpage
\section*{Task 2: Look at what actually happened}

\textbf{Why.} Before trying any fix we wanted the 104 episodes to tell us where the leverage is, without assuming the answer. One rule for this task: describe, don't intervene.

\textbf{What we did.} An agent on a RunPod pod turned every episode into a timeline from the monitor logs (first contact, lid closed, lift, tip, violation), tiled frames from the videos at those moments for all 26 tasks, compared the two same-seed repeats task by task, and put successes next to failures of the same task to see where they first differ. It also wrote a per-step recorder for actions, joint state and jar pose, which the next GPU rental will switch on. The one simulator episode planned never ran: the pod's driver (595) crashes Isaac Sim's renderer at startup, so simulation stays on the 580-driver Vast hosts. Cost about \$1.50.

\textbf{What we found.}
\begin{itemize}\setlength\itemsep{1pt}
  \item The approach is the same every time. Repeats of one seed reach the jar within half a second of each other and look identical until the fingers meet the lid. Everything that decides the outcome happens after contact.
  \item Two moments decide it: whether the push gets the lid over, and the regrasp afterwards, when the gripper leaves the lid, descends, and lands on the body or on the rim. Seven of ten same-seed pairs split in that regrasp.
  \item Idling is as common as crashing: 36 episodes end with nothing lifted and no violation, half of them with the lid already closed and the robot parked. The biggest crash story is tipping the jar after closing it, 19 episodes.
  \item Only 8 tasks behave the same way every run (five always tip, two always park, one always succeeds). The other 18 are knife-edge.
\end{itemize}

\textbf{Where that points.} The lever is in execution and recovery after contact, not in the approach: a slower or shallower regrasp, a stop when a finger meets the rim, a retry when the push has stalled. The leverage note ranks three stage 2 experiments, each about 200 episodes per condition. Panels are frozen: development 0, 3, 6, 9, 15, 17; held out 2, 5, 8, 14, 16, 23, 7, 18, 12.

\textbf{Where it lives.} Repository, under \texttt{progress/receipts/}: the frame sheets (\texttt{02-frames/}), the pair strips (\texttt{02-pairs/}), the timelines, sensitivity map and task notes (\texttt{02-*.md}); the same sheets are in the dataset under \texttt{analysis/02/}. The leverage note, \emph{Where the leverage is}, sits with the other research notes under \texttt{docs/research/}.
\end{document}
